% !TeX root = ../Book1.tex
\begin{figure}[htbp]
	\centering
	\begin{tikzpicture}[x=1.45cm,y=1.25cm,>=stealth,font=\small]
		\draw[->,gray] (-1.7,0)--(1.8,0) node[right]{\(x\)};
		\draw[->,gray] (0,-1.25)--(0,1.35) node[above]{\(y\)};
		\draw[thick,domain=-1.05:1.05,samples=80] plot(\x,{\x*\x*\x});
		\foreach \t in {-0.65,0,0.65}{
			\pgfmathsetmacro{\v}{sqrt(1+9*\t*\t*\t*\t)}
			\draw[blue!70!black,thick,->] (\t,{\t*\t*\t})--({\t-0.7*3*\t*\t/\v},{\t*\t*\t+0.7/\v});
		}
		\node[blue!70!black] at (-1.4,-0.1){\(JT\)};
		\draw[red!70!black,thick,dashed,->] (-0.65,-0.274625)--(-0.0998,-0.7087);
		\node[red!70!black] at (-0.12,-0.95){左侧的 \(N=-JT\)};
		\node[blue!70!black,align=left] at (2.8,0.8){右侧 \(N=JT\)\\但 \(JT\) 始终连续};
		\fill (0,0) circle(1.4pt);
	\end{tikzpicture}
	\caption[拐点两侧的法向选择]{曲线 \(y=x^3\) 的左法向 \(JT\) 连续穿过原点。Frenet 主法向指向弯曲的一侧，因而在左侧取 \(-JT\)、右侧取 \(JT\)，不能在原点连续拼接。}
	\label{g1:fig:inflection-frame}
\end{figure}
